%
% Risch-Algorithmus für algebraische Erweiterungen
%
\subsection{Risch-Algorithmus für algebraische Erweiterungen
\label{buch:intalg:risch:algebraisch}}
Die Methoden, die für logarithmische und exponentielle Erweiterungen
konstruiert wurden, lassen sich nicht direkt auf algebraische
Erweiterungen übertragen.
Dies deutet sich bereits im Beweis des Liouville-Prinzips für algebraische
Erweiterungen an.
Dort wurde zum Beispiel das Konzept der konjugierten Elemente und andere
rein algebraische Eigenschaften verwendet.

Die bisher in Abschnitt~\ref{buch:intalg:section:risch} verwendeten
Methoden basieren wesentlich darauf, dass sich Polynome in $\vartheta$
auf eindeutige Art faktorisieren lassen.
Der euklidische Algorithmus ist für alle Argumente zentral.
Für eine algebraische Erweiterung gilt dies nicht mehr.

Abschnitt~\ref{buch:intalg:section:wurzel} hat aber gezeigt, dass
eine ähnliche Vorgehensweise denkbar ist, um ein Stammfunktion
einer Funktion in $K(x,y)$ zu finden, wobei $y$ ein über $K(x)$
algebraisches Element ist.
Im einfachsten Fall erfüllt $y$ die Gleichung
\[
y^n - h(x) = 0
\]
mit $h(x) \in K(x)$, $y$ ist also eine $n$-te Wurzel einer rationalen
Funktion.

Ein Beispiel sind die in Kapitel~\ref{chapter:elliptisch} diskutierten
\emph{elliptischen Integrale}
\index{elliptisches Integral}%
\[
f(x) = \int_0^x R(t,\!\sqrt{P(t)})\,dt
\]
in denen $p(t)$ ein Polynom vom Grad 3 oder 4 ist.
Der Fall eines Polynoms vom Grad 2 wurde bereits im
Abschnitt~\ref{buch:intalg:section:wurzel} behandelt.

Zunächst muss man eine gute ``Standarddarstellung'' für eine Funktion
$f\in K(x,y)$ finden.
Bei transzendenten Erweiterungen war die Darstellung als gekürzter Bruch
von Polynomen in $\vartheta$ eindeutig, dies gilt für algebraische
Funktionen nicht mehr.
Man wird also zusätzliche Forderungen an Zähler und Nenner aufstellen
müssen, um eine eindeutige Darstellung zu bekommen.

Für die elliptischen Integrale kann man zum Beispiel zeigen, dass sich
jedes elliptische Integral auf die grundlegenden elliptischen Integrale
\begin{align*}
\text{1.~Art:}&&
F(x;k)
&=
\int_0^{\sin x} \frac{dt}{\!\sqrt{(1-t^2)(1-k^2t^2}}
\\
\text{2.~Art:}&&
E(x;k)
&=
\int_0^{\sin x} \frac{\!\sqrt{1-k^2t^2}}{\!\sqrt{1-t^2}}\,dt
\\[5pt]
\text{3.~Art:}&&
\Pi(n,x;k)
&=
\int_0^{\sin x}\frac{dt}{(1-nt^2)\!\sqrt{(1-t^2)}\!\sqrt{1-k^2t^2}}
\end{align*}
zurückführen lässt.

Des Weiteren muss ein Reduktionsverfahren ähnlich dem Verfahren von
Hermite gefunden werden, welches Integrale von rationalen Funktionen
in $K(x,y)$ auf solche von Brüchen mit besonders einfachen Nennern
reduziert.
Die Schwierigkeit dabei ist, dass der euklidische Algorithmus, der
in der Methode von Hermite von zentraler Bedeutung war, nicht mehr
zur Verfügung steht.

Schliesslich müssen die Koeffizienten der logarithmischen Terme
in der Liouville-Form eines Integrals gefunden werden.
Dazu ist ein Prinzip ähnlich der Methode von Rothstein-Trager
nötig.
Die dazu notwendigen Methoden basieren auf der Idee eines Divisors
aus der algebraischen Geometrie.
Damit würden wir aber den Bereich der Analysis auf so grundlegende
Weise verlassen, dass wir uns mit einem Verweis auf die Literatur
begnügen.
\cite{buch:algorithms} enthält eine kurze Skizze der Problem, die 
zu lösen sind und verweist auf weitere relevante Spezialliteratur.








